% !TEX root = main.tex
\documentclass[main.tex]{subfiles}
\begin{document}

\section{Bordered Newton Step}
\label{sec:appendix-bordered}

We detail next the Newton step of Algorithm~\ref{alg:flow-constrained-newton}, which solves \eqref{eq:flow-line-constraint} by maximizing $\jointdual$ subject to the linear constraints \eqref{eq:flow-line-constraint-dil} and \eqref{eq:flow-line-constraint-trans}.

First, we stack the shifts and weights into $\smash{\mathbf{z} \!=\! (\shift, \weight)}$ and the constraint multipliers into $\smash{\boldsymbol{\nu} = 2\,(s, t_1, \dots, t_d) \in \R^{d+1}}$.
We also write the dilation constraint \eqref{eq:flow-line-constraint-dil} as the row $\smash{\mathbf{a} = (\{\pvol_i\,\dual^k_i\}_i, 0)}$ and the translation constraint \eqref{eq:flow-line-constraint-trans} along each axis $\mathbf{e}_j$ as the row $\smash{\mathbf{b}_j = (\{\pvol_i\,\mathbf{e}_j\}_i, 0)}$, both with zero entries on the weights.
We then collect these $d \!+\! 1$ rows into the auxiliary matrix
\begin{equation}
  \Cmat = \begin{bmatrix} \mathbf{a} & \mathbf{b}_1 & \cdots & \mathbf{b}_d \end{bmatrix}\tr ,
  \label{eq:appendix-constraint-matrix}
\end{equation}
so that the constraints \eqref{eq:flow-line-constraint-dil} and \eqref{eq:flow-line-constraint-trans} read $\Cmat\,\mathbf{z} = 0$.
% 
Additionally, we denote the gradient \eqref{eq:joint-gradient} by $\smash{\mathbf{g} \!=\! \nabla\jointdual(\mathbf{z})}$ and the matrix
\begin{equation}
  \Nmat = -\nabla^2\jointdual(\mathbf{z}) + \hreg\,\mathrm{diag}(\pvol),
  \label{eq:appendix-bordered-matrix}
\end{equation}
which is the negated Hessian of $\jointdual$ regularized by $\hreg\,\mathrm{diag}(\pvol)$ with a small $\hreg \!>\! 0$ to remove its null-space.
Since the Hessian is negative semi-definite, $\Nmat$ is positive definite.
Moreover, $\Nmat$ shares the same sparsity as the Hessian, so it admits a sparse Cholesky factorization.
% with the translation $\smash{\translation = (t_1, \dots, t_d)}$, and the Newton step $\delta\mathbf{z}$, 

With this notation, we then solve the bordered system
\begin{equation}
  \begin{bmatrix} \Nmat & \Cmat\tr \\ \Cmat & 0 \end{bmatrix}
  \begin{bmatrix} \delta\mathbf{z} \\ \boldsymbol{\nu} \end{bmatrix}
  =
  \begin{bmatrix} \mathbf{g} \\ 0 \end{bmatrix}.
  \label{eq:flow-constrained-newton}
\end{equation}
%
Eliminating $\delta\mathbf{z}$ from the first row of \eqref{eq:flow-constrained-newton} gives $\smash{\delta\mathbf{z} = \Nmat^{-1} \bigl( \mathbf{g} - \Cmat\tr \boldsymbol{\nu} \bigr)}$, which we compute from the $d \!+\! 2$ sparse solves
\begin{equation}
  \label{eq:flow-constrained-elimination-solve}
  \left\{
  \begin{aligned}
    \Nmat\,\mathbf{z}_0 &= \mathbf{g} , \\
    \Nmat\,\mathbf{z}_a &= 2\,\mathbf{a} , \\
    \Nmat\,\mathbf{z}_{b_j} &= 2\,\mathbf{b}_j \ \ (j = 1, \dots, d) ,
  \end{aligned}
  \right.
\end{equation}
which share a single factorization of $\Nmat$.
Substituting into the last row $\Cmat\,\delta\mathbf{z} = 0$ yields the dense $(d \!+\! 1) \times (d \!+\! 1)$ system
\begin{equation}
    \begin{bmatrix}
      \mathbf{a}\tr\mathbf{z}_a & \mathbf{a}\tr\mathbf{z}_{b_1} & \cdots & \mathbf{a}\tr\mathbf{z}_{b_d} \\
      \mathbf{b}_1\tr\mathbf{z}_a & \mathbf{b}_1\tr\mathbf{z}_{b_1} & \cdots & \mathbf{b}_1\tr\mathbf{z}_{b_d} \\
      \vdots & \vdots & \ddots & \vdots \\
      \mathbf{b}_d\tr\mathbf{z}_a & \mathbf{b}_d\tr\mathbf{z}_{b_1} & \cdots & \mathbf{b}_d\tr\mathbf{z}_{b_d}
    \end{bmatrix}
    \begin{bmatrix} s \\ t_1 \\ \vdots \\ t_d \end{bmatrix}
    =
    \begin{bmatrix} \mathbf{a}\tr\mathbf{z}_0 \\ \mathbf{b}_1\tr\mathbf{z}_0 \\ \vdots \\ \mathbf{b}_d\tr\mathbf{z}_0 \end{bmatrix} ,
    \label{eq:flow-constrained-elimination-gauge}
\end{equation}
which is symmetric positive definite.
The Newton step follows as
\begin{equation}
  \delta\mathbf{z} = \mathbf{z}_0 - s\,\mathbf{z}_a - \sum\nolimits_j t_j\,\mathbf{z}_{b_j}.
  \label{eq:flow-constrained-elimination-update}
\end{equation}

\end{document}
